% App G: Perturbative-regime bound in the torsionful background
% Forward-reference targets:
%   \cref{PerturbativeRegime} from theory.tex §2.2 and §2.3
%   \cref{SurveyCoefficient}  from theory.tex §2.3 EM-conventions paragraph

\section{Perturbative-regime bound in the torsionful background}
\label{PerturbativeRegime}

\paragraph*{Background-perturbativity in Einstein--Maxwell}

Following the standard Gertsenshtein
analysis~\cite{gertsenshtein1962wave,boccaletti1970conversion,raffelt1988mixing},
the linearised calculation takes the magnetic field $\Bzero$ as
externally prescribed rather than self-consistently sourcing the
background metric.
Hwang \& Noh~\cite{hwangnoh2023graviton} examine the validity of
this assumption: the EM stress-energy modifies the linearised
graviton equation of motion through a tachyonic mass-like term
$-(\kappa^2/2)\Bzero^2 h$, with critical scale
$\lambda_B = 2\pi/k_B$ and $k_B = \sqrt{8\pi G}\,\Bzero/c^2$,
above which the dimensionless backreaction parameter
$D^2/\lambda_B^2$ ceases to be small ---
signalling that the standard Gertsenshtein assumption (the magnetic
field's own gravitational field is negligible) breaks down.
They show that the backreaction parameter is proportional to
$\Pconv$:
\begin{align}
  \frac{D^2}{\lambda_B^2}
  \;\sim\; \frac{\Pconv}{\sin^2\!\theta} ,
  \label{HwangNohBound}
\end{align}
where $\theta$ is the angle between $\Bzero$ and the propagation
direction (only the transverse component $\Bzero\sin\theta$
couples)~\cite{hwangnoh2023graviton}.
The conversion probability and the backreaction therefore become
negligible together as $\Bzero\to 0$.

Hwang \& Noh also obtain the peak-amplitude estimate
$P\sim 8\pi^2\sin^2\!\theta\,(D/\lambda_B)^2$ from the coupled
graviton--photon wave equations.
This differs by a factor of four from the small-angle expansion of
\cref{GertsenshteinFormula},
$\Pconv\approx\kappa^2\Bzero^2 D^2\sin^2\!\theta/4$:
Hwang \& Noh estimate the energy ratio
$P\equiv(c^2/16\pi G)|\dot{h}|^2/|b|^2$ by scaling the peak amplitude,
while \cref{GertsenshteinFormula} gives the exact conversion
probability.
The factor of four reflects the difference between linear
build-up and oscillatory mixing.
The exact two-mode oscillation \cref{GertsenshteinFormula} has
graviton amplitude $h\propto\sin(\Omega D/2)$ with mixing frequency
$\Omega = \kappa\Bzero\sin\theta$, the factor of $1/2$ in the argument
being the standard structure of two-mode
mixing~\cite{raffelt1988mixing}.
Hwang \& Noh's peak-amplitude estimate replaces the oscillation by
the linear build-up $h\sim\Omega D$ that the source would produce
in the absence of back-coupling, dropping this $1/2$.
The probability $P\sim|h|^2$ then differs by a factor of four in
the small-$\Omega D$ regime.
\cref{GertsenshteinFormula} is the authoritative quantitative form;
\cref{HwangNohBound} is used as a parametric bound only.
Equivalently, $\kappa\Bzero D\ll 1$, i.e. $\Pconv\ll 1$.

\paragraph*{Self-consistency at vanishing background field}

PGT coupled to electromagnetism adds a torsion perturbation
$\delta\MAGT{^\sigma_\mu_\nu}$ alongside the graviton $h_{\mu\nu}$
and photon perturbation $\delta\FieldA{_\mu}$, and non-minimal
couplings such as $\Del{1}\MAGFt{_{[\mu\nu]}}\FieldF{^{\mu\nu}}$
introduce new mixing vertices.
At $\Bzero=0$ the configuration $g_{\mu\nu}=\eta_{\mu\nu}$,
$\bar{A}_\mu=0$, $\bar{\MAGT{}}=0$ ---
the background of \cref{Linearisation} evaluated at $\Bzero=0$ ---
satisfies the full nonlinear field equations of any ghost-safe
polynomial PGT Lagrangian identically: every field-equation term
is at least linear in the field strengths
$\{\widetilde{\mathcal{R}},\MAGT{},\FieldF{},\rcD{}\MAGT{}\}$,
all of which vanish on this background.
Each such term therefore vanishes term-by-term, and their sum ---
the field equation --- vanishes identically.
The leading Taylor coefficient $c_2$ of the expansion
\cref{TaylorPconv} is therefore evaluated free of the
metric-backreaction shift discussed above.

\paragraph*{Analytic dependence on the background field}

The background field $\bar{\FieldF{}}_{zy}=-\Bzero$ enters the
linearised equations of motion only through coupling vertices,
so the kinetic matrix $\mathcal{K}$ of \cref{KineticMatrix}
acquires an explicit $\Bzero$-dependence
$\mathcal{K}(\partial_t,\partial_z;\Bzero)$ that is polynomial in
$\Bzero$ (at most linear for a quadratic Lagrangian, at most
quadratic if cubic terms are included).
The conversion amplitude follows from $\mathcal{K}$ by standard
mode decomposition, yielding a transition amplitude that is an
entire function of $\Bzero$.
At $\Bzero=0$ every graviton--photon coupling vertex carries at
least one contraction with $\bar{\FieldF{}}\propto\Bzero$;
all such vertices vanish, the kinetic matrix is block-diagonal in
the graviton--photon channel, and $\Pconv(0)=0$.
The Taylor expansion in $\Bzero$ is therefore
\begin{align}
  \Pconv \;=\; c_2\Bzero^2 + c_4\Bzero^4 + \cdots ,
  \label{TaylorPconv}
\end{align}
with $c_2=\Czero$ finite.

\paragraph*{Vanishing of the background torsion}

On the uniform-field flat background
$\partial_\mu\bar{\FieldF{}}_{\nu\rho}=0$.
Varying $\Del{1}\int\bar{\FieldF{}}^{\mu\nu}\delta\widetilde{\mathcal{R}}_{[\mu\nu]}$
with respect to $\MAGT{}^{\rho}{}_{\mu\nu}$ and integrating by parts
gives $\Del{1}(\nabla_\rho\bar{\FieldF{}}^{\mu\nu})\cdot(\,\cdot\,)=0$.
The remaining Cartan-equation terms vanish because
$\widetilde{\mathcal{R}}_{[\mu\nu]}|_{\eta,\,T=0}=0$
(the antisymmetric Ricci of flat Minkowski with zero torsion vanishes
identically) and all torsion mass and kinetic terms are at least
linear in $\MAGT{}$.
Thus $\bar{\MAGT{}}=0$ is an exact solution of the Cartan equation
on this background, not merely an approximation.
This conclusion relies on $\partial_\mu\bar{\FieldF{}}=0$:
for non-uniform fields or curved backgrounds
$\nabla\bar{\FieldF{}}\neq 0$ and $\bar{\MAGT{}}=0$ need not
hold~\cite{bahamonde2025torsion}.

\paragraph*{Torsion-perturbation bound}

The linearised Cartan equation for $\delta\MAGT{}$ has its
leading-order $\bar{\FieldF{}}\nabla\delta\MAGT{}$ coupling cancelled
by the same integration by parts.
The source comes from the contribution to
$\widetilde{\mathcal{R}}_{[\mu\nu]}$ that is bilinear in $h$ and
$\delta\MAGT{}$ --- schematically $(\partial h)\cdot K[\delta\MAGT{}]$,
with $K$ the contortion built from $\delta\MAGT{}$.
Varying with respect to $\delta\MAGT{}$ produces a source
$\propto\Del{1}\bar{\FieldF{}}\,\partial h$, of order
$\Del{1}\Bzero A/D$ for a plane-wave graviton of dimensionless
amplitude $A$ (the IC amplitude of \cref{Linearisation})
propagating over distance $D$.
Algebraic inversion through the $\Alp{}\MAGT{}^2$ mass term gives
$|\delta\MAGT{}|\sim(|\Del{1}|/\Alp{})\,\Bzero A$, with the wavenumber
absorbed at the simulation scale.
Perturbativity of the torsion sector
($|\delta\MAGT{}|\ll\Alp{}$) requires
\begin{align}
  |\Del{1}|\,\Bzero\,A \;\ll\; \Alp{}^{\,2} .
  \label{PerturbativeBound}
\end{align}
At the worst-case point of the survey
$(\Bzero,A,|\Del{1}|,\Alp{}) = (10^{-3},0.1,1,1)$
this reads $10^{-4}\ll 1$;
at the smallest stable $\Alp{}\simeq 0.1$, still $10^{-4}\ll 10^{-2}$.

\paragraph*{Generality beyond the algebraic-torsion theory}

The preceding properties hold for any ghost-safe polynomial PGT
Lagrangian coupled to electromagnetism~\cite{blagojevic2018hamiltonian}.
The bound \cref{PerturbativeBound} is specific to the algebraic-torsion
theory with the
$\Del{1}\MAGFt{_{[\mu\nu]}}\FieldF{^{\mu\nu}}$ coupling.
For a generic non-minimal coupling of schematic form
$g_{\mathrm{NM}}\,\widetilde{\mathcal{R}}\,\FieldF{}$ ---
with $\Del{1}$ the specific instance considered above ---
and a torsion mass scale $m_T$ set by the algebraic torsion term
($m_T^2\sim\Alp{}$ for the campaign theory),
the torsion correction to the survey coefficient is
$O\!\left(\Bzero^2(g_{\mathrm{NM}}/m_T^2)^2\right)$.
For theories without non-minimal coupling the torsion sector
decouples at $\Bzero=0$ and contributes nothing.

\paragraph*{Self-consistency of the sector survey}

Across the parameter regions of \cref{SectorSurvey} ---
$\Bzero\in[10^{-5},10^{-3}]$ in simulation units,
$A=0.1$, with coupling coefficients within their swept windows ---
\cref{PerturbativeBound} is satisfied with at least two decades of
margin and $\Pmax<0.1$ throughout,
confirming \cref{HwangNohBound}.
\cref{PerturbativityHierarchy} collects the four linearisation
conditions and their dependence on $\Bzero$.

\begin{table}[htbp]
  \centering
  \caption{Linearisation conditions for the linearised Gertsenshtein
    calculation in the torsionful background and their dependence
    on $\Bzero$.
    The first and fourth conditions are set by the IC and geometry
    choices of \cref{Linearisation};
    the middle two are satisfied as $\Bzero\to 0$
    at fixed IC amplitude~$A$.}
  \label{PerturbativityHierarchy}
  \begin{tabular}{@{}lll@{}}
    \toprule
    Condition & Physical origin & $\Bzero$-dependence \\
    \midrule
    $\lvert h_{\mu\nu}\rvert \ll 1$
      & metric linearisation
      & independent of $\Bzero$ \\
    $\Pconv \ll 1$
      & Hwang--Noh bound \cref{HwangNohBound}
      & $\propto\Bzero^2$ \\
    $\lvert\delta\MAGT{}\rvert \ll \Alp{}$
      & torsion-sector linearisation \cref{PerturbativeBound}
      & $\propto\Bzero$ \\
    $\lvert\delta\FieldA{_\mu}\rvert \ll \lvert\bar{A}_\mu\rvert$
      & EM linearisation
      & independent of $\Bzero$ \\
    \bottomrule
  \end{tabular}
\end{table}

The multi-$\Bzero$ convergence diagnostic at
$\Bzero\in\{10^{-3},10^{-4},10^{-5}\}$ is the perturbative-regime check:
the $c_4\Bzero^2$ remainder in \cref{TaylorPconv} predicts
$C(\Bzero)\to\Czero$ at a rate linear in $\Bzero^2$, and this is what
the survey of \cref{SectorSurvey} measures.

\paragraph*{The survey coefficient}

The parameter survey of \cref{SectorSurvey} evaluates the limiting
survey coefficient
\begin{align}
  \Czero \;\equiv\; \lim_{\Bzero\to 0}\,\frac{\Pconv}{\Bzero^2}
  \;=\; c_2 ,
  \label{SurveyCoefficient}
\end{align}
the leading Taylor coefficient of \cref{TaylorPconv}.

All corrections to $C = \Pconv/\Bzero^2$ enter at $O(\Bzero^2)$,
leaving $\Czero$ at $O(1)$.
The Taylor remainder is $c_4\Bzero^2$ from \cref{TaylorPconv}.
The metric backreaction enters via the tachyonic-mass term
$-(\kappa^2/2)\Bzero^2 h$ of the Einstein--Maxwell paragraph above,
which detunes the graviton dispersion by $\Delta\propto\Bzero^2$.
The mixing of two modes coupled at strength $g\propto\Bzero$ with
detuning $\Delta\propto\Bzero^2$ yields the standard two-mode
oscillation~\cite{raffelt1988mixing}
\begin{align*}
  \Pconv \;=\; \frac{g^2}{g^2+\Delta^2}
  \,\sin^2\!\!\left(\!\sqrt{g^2+\Delta^2}\,D\right) ;
\end{align*}
substituting $\Delta/g\propto\Bzero$ gives a correction to $C$
at $O(\Bzero^2)$.
The torsion nonlinearity contributes at
$O(\Bzero^2(g_{\mathrm{NM}}/m_T^2)^2)$ per
\cref{PerturbativeBound}.
Signal and corrections are separated by one power of $\Bzero^2$.

The amplification factor
$\Aamp\equiv\sqrt{\Czero^{\mathrm{torsion}}/\Czero^{\mathrm{EM}}}$
of \cref{SectorSurvey} is $\Bzero$-independent to $O(\Bzero^2)$
by the same argument.
The limiting coefficient $\Czero^{\mathrm{EM}}$ matches the
potential-formulation derivation of~\cite{hwangnoh2023graviton}:
their conversion coefficient $\Delta_{g\gamma}=\kappa\Bzero/2$
is exactly linear in $\Bzero$, so $\Czero^{\mathrm{EM}}$ has no
$\Bzero$-independent backreaction shift.
